\notepage{35}{Classically Allowed and Forbidden Regions}
Classically the kinetic energy is $E-V(x)$. Motion is allowed where
$E\ge V(x)$ and forbidden where $E<V(x)$. The source illustrates smooth
and square wells, infinite walls, and a well approaching a finite value:
\begin{center}
\begin{tikzpicture}[x=0.65cm,y=0.6cm,>=Stealth]
 \begin{scope}
  \draw[->] (-1.6,0)--(1.6,0) node[right] {$x$};
  \draw[ink,thick,domain=-1.45:1.45] plot (\x,{0.7*\x*\x});
  \draw[dashed] (-1.5,0.8)--(1.5,0.8) node[right] {$E$};
 \end{scope}
 \begin{scope}[xshift=3.9cm]
  \draw[->] (-1.6,0)--(1.6,0) node[right] {$x$};
  \draw[ink,thick] (-1.5,1.45)--(-0.9,1.45)--(-0.9,0.1)--(0.9,0.1)--(0.9,1.45)--(1.5,1.45);
  \draw[dashed] (-1.5,0.8)--(1.5,0.8) node[right] {$E$};
 \end{scope}
 \begin{scope}[xshift=7.8cm]
  \draw[->] (-1.6,0)--(1.6,0) node[right] {$x$};
  \draw[ink,thick,->] (-1.1,0.2)--(-1.1,1.6);
  \draw[ink,thick,->] (1.1,0.2)--(1.1,1.6);
  \draw[ink,thick] (-1.1,0.2) to[bend right=18] (1.1,0.2);
  \draw[dashed] (-1.5,0.8)--(1.5,0.8) node[right] {$E$};
 \end{scope}
 \begin{scope}[xshift=11.7cm]
  \draw[->] (-1.6,0)--(1.6,0) node[right] {$x$};
  \draw[ink,thick,domain=-1.5:1.5,samples=50] plot (\x,{1.5-1.4*exp(-2*\x*\x)});
  \draw[dashed] (-1.5,0.8)--(1.5,0.8) node[right] {$E$};
 \end{scope}
\end{tikzpicture}
\end{center}
\smallheading{Locally classically allowed: oscillatory solutions}
In a small interval around $x_0$ where the potential is approximately constant,
\[
 \frac{2m}{\hbar^2}[E-V(x)]\simeq k^2\ge0,
 \qquad\psi''\simeq-k^2\psi,
 \qquad\psi\simeq A\ee^{\ii kx}+B\ee^{-\ii kx}.
\]
\smallheading{Locally classically forbidden: exponential solutions}
If $E<V(x)$ and $V$ is approximately constant in the interval,
\[
 \frac{2m}{\hbar^2}[E-V(x)]\simeq-\kappa^2<0,
 \qquad\psi''\simeq\kappa^2\psi,
 \qquad\psi\simeq A\ee^{\kappa x}+B\ee^{-\kappa x},\quad\kappa>0.
\]
\begin{center}
\begin{tikzpicture}[x=0.8cm,y=0.6cm,>=Stealth]
 \begin{scope}
  \draw[->] (-1.7,0)--(1.7,0) node[right] {$x$};
  \draw[ink,domain=-1.5:1.5] plot (\x,{0.4+0.3*\x*\x});
  \draw[dashed] (-1.6,1.3)--(1.6,1.3) node[right] {$E$};
  \draw[dotted] (-0.3,0)--(-0.3,0.43);\draw[dotted] (0.3,0)--(0.3,0.43);
  \node at (0,-0.5) {allowed near $x_0$};
 \end{scope}
 \begin{scope}[xshift=4.5cm]
  \draw[->] (-1.7,0)--(1.7,0) node[right] {$x$};
  \draw[ink] (-1.5,1.5) to[bend right=18] (1.5,0.7) node[right,above] {$V$};
  \draw[dashed] (-1.6,0.5) node[left] {$E$} --(1.6,0.5);
  \node at (0,-0.5) {forbidden};
 \end{scope}
 \begin{scope}[xshift=9cm]
  \draw[->] (-1.7,0)--(1.7,0) node[right] {$x$};
  \draw[ink,domain=-1.4:1.4] plot (\x,{0.7*exp(0.6*\x)});
  \draw[dashed] (-1.6,0.5)--(1.6,0.5) node[right] {$E$};
  \node at (0,-0.5) {forbidden tail};
 \end{scope}
\end{tikzpicture}
\end{center}
If $E<V(x)$ for all $x>x_0$, a physical bound-state wavefunction must
exclude the growing right-tail solution and decay as $x\to+\infty$.
Likewise, a forbidden left tail must decay as $x\to-\infty$.
For a confining potential with $V(x)\to+\infty$ as $|x|\to\infty$,
for each $E$ there may be turning points
$x_1>x_2$ with $E<V(x)$ for $x>x_1$ and $x<x_2$; physical eigenstates
decaying at both ends are bound states.
\editorial{The source writes $A(x)\ee^{\kappa(x)x}+B(x)\ee^{-\kappa(x)x}$.
This is local-constant shorthand, not a general solution for variable
$\kappa$. Where a WKB approximation applies, the tails instead have
the form $\kappa(x)^{-1/2}\exp[\pm\int^x\kappa(s)\,\dd s]$;
turning points need separate treatment.}
